\cleardoublepage
\phantomsection
\addcontentsline{toc}{part}{The Igusa Square Root}
\chapter*{Book entrance: The Igusa Square Root}
\addcontentsline{toc}{chapter}{Book entrance: The Igusa Square Root}
\begin{summarybox}
The Jacobi seed, Borcherds denominator, root superalgebra, local--wrapped currents, quantum group, Hall model, and Fredholm realization.
\end{summarybox}
\part{The automorphic determinant}

\chapter{The protected Jacobi input}
\section{The K3 elliptic genus}
Let $S$ be a K3 surface. The right-moving supersymmetric index of a unitary $\mathcal N=(4,4)$ sigma model is the elliptic genus
\[
 Z_{K3}(\tau,z)=2\phiweak(\tau,z).
\]
The normalized weak Jacobi form has weight zero, index one, and Fourier expansion
\[
 \phiweak(\tau,z)=\sum_{n\ge0}\sum_{\ell\in\Z}f(n,\ell)q^nr^\ell.
\]
In Eichler--Zagier normalization,
\[
 \phiweak(\tau,z)=4\sum_{a=2}^{4}
 \frac{\vartheta_a(\tau,z)^2}{\vartheta_a(\tau,0)^2}
 \in J^{\mathrm w}_{0,1}
\]
\cite{EZ}. The coefficient depends on the discriminant $4n-\ell^2$ and on $\ell$ modulo $2$.

\begin{theorem}[Exact low Fourier coefficients]
The first four $q$-rows are
\[
\begin{aligned}
\phiweak={}&r^{-1}+10+r\\
&+q\bigl(10r^{-2}-64r^{-1}+108-64r+10r^2\bigr)\\
&+q^2\bigl(r^{-3}+108r^{-2}-513r^{-1}+808-513r+108r^2+r^3\bigr)\\
&+q^3\bigl(-64r^{-3}+808r^{-2}-2752r^{-1}+4016\\
&\hspace{7.5em}{}-2752r+808r^2-64r^3\bigr)+O(q^4).
\end{aligned}
\]
In particular,
\[
 f(0,0)=10,
 \qquad \frac{f(0,0)}2=5.
\]
\end{theorem}
\begin{proof}
Insert the theta-series expansions into the preceding quotient formula. Work with $t=q^{1/8}$ and $x=r^{1/2}$; each quotient is an exact formal power series in $t$ with Laurent-polynomial coefficients in $x$. The fractional powers cancel in the sum. Division through $t^{24}$ gives the displayed rows. The arithmetic calculation is reproduced in Appendix~\ref{app:theta-computation}.
\end{proof}

\section{The virtual half-index}
For every $(n,\ell)$ choose a finite-dimensional super vector space $U_{n,\ell}$ satisfying
\[
 \sdim U_{n,\ell}=f(n,\ell).
\]
The integer \(f(n,\ell)\) is its superdimension. In the ordinary
Grothendieck group of finite-dimensional super vector spaces, the
even and odd dimensions remain separate. The signed integer is their
image under superdimension. The determinant depends only on that image.

\begin{lemma}[Fock character and Berezinian]
Let $V=\bigoplus_{\gamma}V_\gamma$ be a locally finite graded super vector space. Then
\[
 \sch\bigl(\Sym(V_{\bar0})\tensor\bigwedge(V_{\bar1})\bigr)
 =\prod_\gamma(1-x^\gamma)^{-\sdim V_\gamma},
\]
and
\[
 \sdet_V(1-x)=\prod_\gamma(1-x^\gamma)^{\sdim V_\gamma}.
\]
\end{lemma}
\begin{proof}
An even one-dimensional summand contributes the geometric series $(1-x^\gamma)^{-1}$. An odd one-dimensional summand contributes $1-x^\gamma$. Multiplication over a homogeneous basis proves both identities.
\end{proof}

\chapter{The Borcherds product}
\section{The effective chamber}
Define
\[
\begin{aligned}
\Gammaeff={}&\set{(n,\ell,m)\in\Z^3:m>0,\ n\ge0}\\
&\cup\set{(n,\ell,0):n>0}
 \cup\set{(0,\ell,0):\ell<0}.
\end{aligned}
\]
The ordering is lexicographic in $(m,n,-\ell)$ and is precisely the ordering selected by the cusp chamber.
For $\gamma=(n,\ell,m)$ put
\[
 V_\gamma=U_{nm,\ell},
 \qquad \sdim V_\gamma=f(nm,\ell).
\]

\begin{definition}[Monic virtual determinant]
The monic determinant of the half-index is
\[
 \Dmon(Z)=q^{1/2}r^{1/2}s^{1/2}
 \sdet_V(1-x)
 =q^{1/2}r^{1/2}s^{1/2}
 \prod_{\Gammaeff}(1-q^nr^\ell s^m)^{f(nm,\ell)}.
\]
\end{definition}

\begin{theorem}[Borcherds--Gritsenko--Nikulin product]
The preceding product converges normally on the chosen cusp chamber and extends to a holomorphic genus-two modular form of weight five with a character of order two. It is the monic Igusa form $\Dmon$. The theta-normalized form is
\[
 \Dtheta=2^6\Dmon.
\]
\end{theorem}
\begin{proof}
The singular theta lift of the vector-valued form corresponding to $\phiweak$ produces the product, divisor, weight, and automorphic character. The weight is $f(0,0)/2=5$. The genus-two identification with the Igusa form and the product formula are the Gritsenko--Nikulin theorems \cite{Borcherds95,Borcherds98,GN97,GN98b}. The product of the ten even theta constants has coefficient $2^6$ at $q^{1/2}r^{1/2}s^{1/2}$, fixing the scalar.
\end{proof}

\section{The theta constant normalization}
Let $\mathfrak E$ be the ten even genus-two characteristics. Define
\[
 \Dtheta(Z)=\prod_{\mathfrak m\in\mathfrak E}\vartheta[\mathfrak m](Z,0).
\]
The cusp leading terms of the ten factors contribute six factors of $2$ and total exponent $(1/2,1/2,1/2)$. Thus
\[
 [q^{1/2}r^{1/2}s^{1/2}]\Dtheta=64.
\]
This coefficient selects a canonical scalar normalization of the automorphic section.

\begin{proposition}[Ten-factor leading-term calculation]
Group the ten even characteristics by the first half-characteristic $a\in(\Z/2)^2$. Their leading contributions in the chosen cusp chamber are
\[
\begin{array}{c|c|c|c}
a&\#&\text{scalar}&\text{monomial}\\ \hline
(0,0)&4&1&1\\
(1,0)&2&2&q^{1/8}\\
(0,1)&2&2&s^{1/8}\\
(1,1)&2&2&q^{1/8}r^{1/4}s^{1/8}.
\end{array}
\]
Multiplication gives
\[
 1^4\,2^2\,2^2\,2^2\,
 q^{1/2}r^{1/2}s^{1/2}
 =64q^{1/2}r^{1/2}s^{1/2}.
\]
\end{proposition}
\begin{proof}
Insert each characteristic into the genus-two theta series and retain the unique terms of least order in the lexicographic cusp valuation $(m,n,-\ell)$. For $a=(0,0)$ the zero lattice vector contributes one. For each nonzero $a$, the two shortest lattice representatives are exchanged by sign and contribute the scalar two. Their quadratic exponents are the entries in the final column. Summing the exponents and multiplying the six factors of two gives the formula.
\end{proof}

\begin{proposition}[Boundary factor of the Borcherds product]
The $m=0$ part of the monic product is
\[
 \prod_{n>0,\ell\in\Z}(1-q^nr^\ell)^{f(0,\ell)}
 \prod_{\ell<0}(1-r^\ell)^{f(0,\ell)}
 =(1-r^{-1})\prod_{n\ge1}
 (1-q^nr^{-1})(1-q^n)^{10}(1-q^nr).
\]
Consequently the first Fourier--Jacobi coefficient is
\[
 q^{1/2}r^{1/2}(1-r^{-1})
 \prod_{n\ge1}(1-q^nr^{-1})(1-q^n)^{10}(1-q^nr)
 =\eta(\tau)^9\vartheta_{11}(\tau,z),
\]
where the odd theta function is fixed by the exact convention
\[
 \vartheta_{11}(\tau,z)
 =-q^{1/8}r^{-1/2}
 \prod_{n\ge1}(1-q^{n-1}r)(1-q^nr^{-1})(1-q^n).
\]
\end{proposition}
\begin{proof}
The $q^0$ row of $\phiweak$ has exactly three nonzero coefficients:
$f(0,-1)=1$, $f(0,0)=10$, and $f(0,1)=1$. Substitution into the $m=0$ factors gives the first formula. The Jacobi triple product and $\eta(\tau)=q^{1/24}\prod_{n\ge1}(1-q^n)$ give the second, including its leading sign and monomial.
\end{proof}

\begin{corollary}[The scalar square]
The monic weight-ten form is
\[
 \ChiTen=\Dmon^2=2^{-12}\Dtheta^2.
\]
Its automorphic character is trivial.
\end{corollary}
\begin{proof}
The identity $\Dtheta=2^6\Dmon$ gives the scalar factor.  The character of $\Dmon$ has order two, so its square is trivial.
\end{proof}

\chapter{Divisor, character, and the meaning of the square root}
The diagonal Humbert divisor is
\[
 H_1=\set{Z\in\HHtwo:z=0}/\operatorname{Sp}_4(\Z).
\]

\begin{theorem}[Diagonal divisor]
The divisor of $\Dmon$ is the $\operatorname{Sp}_4(\Z)$-orbit of $H_1$, with multiplicity one. The form transforms by the Maass character $\nu_5$ of order two, and $\ChiTen=\Dmon^2$ is scalar-valued.
\end{theorem}
\begin{proof}
This is the defining diagonal-divisor property of the Igusa square root and the first row of the Cl\'ery--Gritsenko classification \cite{Igusa62,Igusa64,GC}.
\end{proof}

Real roots of the denominator algebra define Weyl reflections. Imaginary roots contribute root spaces and product exponents. This division of labor is intrinsic to a Borcherds algebra: the Weyl group is generated by real simple roots, while the imaginary simple roots enter the correction term of the denominator identity.

\section{Automorphic section, orientation line, and Pfaffian operator}
An automorphic square root is a section of a line bundle with character. An orientation is a square root of a determinant line. A Pfaffian operator is a skew family whose regularized Pfaffian is a section of that orientation line. These form a sequence of successive refinements:
\[
\begin{gathered}
 \text{automorphic section}
 \longrightarrow \text{determinant-line identification}\\
 \longrightarrow \text{oriented square-root line}
 \longrightarrow \text{Fredholm family and Pfaffian section}.
\end{gathered}
\]

\begin{proposition}[Logical separation of the four levels]
The product theorem constructs the first object in the sequence. A determinant-line comparison constructs the second. A coherent square root of that line constructs the third. A first-order operator requires a symbol, domain, ellipticity or Fredholm condition, and regularization, and thereby produces the fourth refinement.
\end{proposition}
\begin{proof}
Each stage adds a specified geometric structure. A determinant realization adds an isomorphism from a geometric determinant line to the automorphic line. An oriented realization adds a coherent square-root line and its parity refinement. An operator realization adds a principal symbol, a Fredholm domain, a spectrum, and a regularization whose Pfaffian is the resulting section.
\end{proof}

The title names the monic automorphic section $\Dmon$, characterized by the exact identity $\Dmon^2=\ChiTen$. Part~IV formulates an analytic Pfaffian realization as a further geometric refinement.

\part{The denominator algebra}

\chapter{The Lorentzian root lattice}
Let
\[
 L=\Z f_2\oplus\Z f_3\oplus\Z f_{-2}
\]
with bilinear form
\[
 (f_2,f_{-2})=-1,
 \qquad (f_3,f_3)=2,
\]
and all remaining basis pairings zero. Its signature is $(2,1)$. Define the type-II root lattice
\[
 L_{II}=2\Z f_2\oplus\Z f_3\oplus2\Z f_{-2}
 \cong U(4)\oplus\angles{2}.
\]
The Gram-index map is
\[
 \alpha:\Gammaind=\Z^3\longrightarrow L_{II},
 \qquad
 \alpha(n,\ell,m)=2nf_2-\ell f_3+2mf_{-2}.
\]
Then
\[
 (\alpha(n,\ell,m),\alpha(n',\ell',m'))
 =-2\bigl(2nm'+2n'm-\ell\ell'\bigr),
\]
and
\[
 (\alpha(n,\ell,m),\alpha(n,\ell,m))
 =-2(4nm-\ell^2).
\]

\section{Real simple roots and Weyl vector}
Put
\[
 \delta_1=2f_2-f_3,
 \qquad
 \delta_2=2f_{-2}-f_3,
 \qquad
 \delta_3=f_3.
\]

\begin{proposition}[Type-II Cartan matrix]
The Gram matrix of the real simple roots is
\[
 \bigl((\delta_i,\delta_j)\bigr)_{i,j=1}^3
 =\begin{pmatrix}
 2&-2&-2\\
 -2&2&-2\\
 -2&-2&2
 \end{pmatrix}.
\]
The reflections
\[
 s_i(x)=x-(x,\delta_i)\delta_i
\]
are integral involutions preserving $L_{II}$ and its form.
\end{proposition}
\begin{proof}
Substitution into the basis pairing gives the matrix. Since $(\delta_i,\delta_i)=2$, the reflection formula has integral coefficients. Direct expansion gives $s_i^2=1$ and $(s_ix,s_iy)=(x,y)$.
\end{proof}

The Weyl vector is
\[
 \rho=f_2-\frac12f_3+f_{-2},
 \qquad (\rho,\delta_i)=-1.
\]
In the Gram-index variables, $e^{-\rho}=q^{1/2}r^{1/2}s^{1/2}$.

\chapter[The denominator algebra]{The Gritsenko--Nikulin Lie superalgebra}
The Gritsenko--Nikulin construction specifies a generalized
Kac--Moody Lie superalgebra \(\gIg\), with real simple roots
\(\delta_1,\delta_2,\delta_3\) and the imaginary correction
attached to \(\phiweak\). Its generators and relations supply the
bracket whose denominator is compared with the Borcherds product.

\begin{theorem}[Root supermultiplicities and denominator]
On the active positive support,
\[
 \smult\gIg_{\alpha(n,\ell,m)}=f(nm,\ell).
\]
In the Lorentzian coordinate convention of Gritsenko--Nikulin,
\[
 \den(\gIg)=\Dmon(2Z)=2^{-6}\Dtheta(2Z).
\]
\end{theorem}
\begin{proof}
The real roots and Weyl chamber are the type-II case of the Gritsenko--Nikulin construction. The imaginary simple-root correction is chosen so that its Weyl--Kac--Borcherds product has exponent $f(nm,\ell)$. The resulting product is the doubled-variable monic Igusa form \cite{GN97,GN98a,GN98b}.
\end{proof}

The theorem gives a Lie superalgebra with a bracket, Cartan, Weyl group, root parity, and denominator. The scalar product alone records the signed dimensions of the root spaces; the Gritsenko--Nikulin construction supplies the algebraic operations.

\chapter{The ordered bar as the denominator complex}
Let $\nIg$ be the positive nilpotent subalgebra of $\gIg$. Fix a proper integral height $h$ on the positive root cone and define
\[
 F^{>H}\nIg=\bigoplus_{h(\alpha)>H}\gIg_\alpha,
 \qquad
 \nIg^{(H)}=\nIg/F^{>H}\nIg.
\]
Height additivity makes $F^{>H}\nIg$ an ideal. Properness and local finiteness make $\nIg^{(H)}$ finite-dimensional and nilpotent.

\begin{theorem}[Height-complete Bar--Chevalley comparison]
For every $H$, antisymmetrization gives a natural quasi-isomorphism of conilpotent dg coalgebras
\[
 \Barc\bigl(U(\nIg^{(H)})\bigr)\simeq C_*(\nIg^{(H)}).
\]
The quotient maps preserve these comparisons and induce a quasi-isomorphism in the separated height-complete category
\[
 \Barc_h^\wedge\bigl(U(\nIg)\bigr)
 :=\varprojlim_H\Barc\bigl(U(\nIg^{(H)})\bigr)
 \simeq
 C_{*,h}^\wedge(\nIg)
 :=\varprojlim_H C_*(\nIg^{(H)}).
\]
Its root-graded Euler supercharacter is
\[
 \chi_{\mathrm{root}}\,C_{*,h}^\wedge(\nIg)
 =\prod_{\alpha>0}
 (1-e^{-\alpha})^{\sdim\gIg_\alpha}.
\]
\end{theorem}
\begin{proof}
For each $H$, the antisymmetrization map sends $\Sym(\nIg^{(H)}[1])$ to the normalized bar of $U(\nIg^{(H)})$. The PBW filtrations have the same associated-graded Koszul resolution of $\Sym(\nIg^{(H)})$, so antisymmetrization is a filtered quasi-isomorphism. Naturality gives a morphism of inverse systems. Every fixed root degree is visible in all sufficiently large height quotients and is unchanged thereafter; hence the inverse-limit map is a quasi-isomorphism coefficientwise. An even root vector becomes odd after the Chevalley shift and contributes $1-e^{-\alpha}$; an odd root vector becomes even and contributes its reciprocal. Multiplication over the finite-dimensional root spaces gives the product.
\end{proof}

Let $\mathbb T_{L_{II}}$ be the formal root torus and write $Z_{\mathrm{root}}$ for its logarithmic coordinate, normalized by
\[
 e^{-\alpha(n,\ell,m)}=Q^nR^\ell S^m,
 \qquad
 e^{-\rho}=Q^{1/2}R^{1/2}S^{1/2}.
\]

\begin{corollary}[Bar denominator identity]
On the formal root torus,
\[
 e^{-\rho}\chi_{\mathrm{root}}
 C_{*,h}^\wedge(\nIg)
 =\Dmon(Z_{\mathrm{root}}).
\]
The Gritsenko--Nikulin embedding of the root torus into the Siegel period domain sends $Z_{\mathrm{root}}$ to $2Z$ \cite{GN97,GN98b}. Hence
\[
 e^{-\rho}\chi_{\mathrm{root}}
 C_{*,h}^\wedge(\nIg)
 =\den(\gIg)=\Dmon(2Z)
\]
in Siegel period coordinates.
\end{corollary}
\begin{proof}
Multiply the Euler product in the preceding theorem by the Weyl monomial $e^{-\rho}$.  The denominator theorem identifies the result on the root torus, and the stated embedding gives the Siegel-coordinate formula.
\end{proof}

The denominator is therefore the Weyl-shifted Euler supercharacter of the height-complete interval bar of the positive enveloping algebra.

\chapter{Finite-height currents}
Retain the finite-height quotients $\nIg^{(H)}$ of the preceding chapter.  Define the current Lie conformal superalgebra
\[
 \Cur(\nIg^{(H)})=\kfield[\partial]\tensor\nIg^{(H)},
 \qquad [a_\lambda b]=[a,b],
\]
and its universal chiral enveloping algebra
\[
 V_H=U^{\ch}\!\left(\Cur_E(\nIg^{(H)})\right).
\]

\begin{theorem}[Height-complete current envelope]
The quotient maps induce a coefficientwise inverse limit
\[
 \widehat V_{\Delta_5}=\varprojlim_HV_H.
\]
Every $V_H$ is PBW-admissible as an associative chiral algebra, the current generators are primitive for the standard current coproduct, and
\[
 \sch\widehat V_{\Delta_5}
 =\prod_{j\ge1}\prod_{\alpha>0}
 (1-t^je^{-\alpha})^{-\sdim\gIg_\alpha}.
\]
Moreover,
\[
 \BarAssCh(V_H)
 \simeq
 \BarLieCh\!\left(\Cur_E(\nIg^{(H)})\right),
\]
and these equivalences pass to the separated height completion.
\end{theorem}
\begin{proof}
PBW orders the negative modes $a(-j)$, $j\ge1$.  Even modes contribute symmetric powers and odd modes contribute exterior powers, giving the product.  Proper height makes every coefficient stationary.  The current coproduct sends $a(z)$ to $a(z)\tensor1+1\tensor a(z)$.  Chiral antisymmetrization identifies the associative-chiral bar of the chiral enveloping algebra with the chiral Chevalley coalgebra; the complete separated PBW filtration is coefficientwise finite.
\end{proof}

\section{Classical local theory and exact quantum zero modes}
For every $H$, the finite-dimensional nilpotent Lie superalgebra $\nIg^{(H)}$ defines a cotangent BF theory on $E\times\R$.  Its local dg Lie algebra is
\[
 \Lcal_H=\Omega^{0,*}(E)\widehattensor\Omega^*(\R)\tensor\nIg^{(H)}
\]
and the classical action is
\[
 S_H(A,B)=\int_{E\times\R}\Omega_E\wedge
 \left\langle B,(\dbar+\dd_{\R})A+\frac12[A,A]\right\rangle .
\]

\begin{theorem}[Finite-height BF realization]
The action $S_H$ satisfies the classical master equation and defines a classical holomorphic--topological factorization algebra.  Harmonic reduction along a compact interval produces the finite-dimensional cotangent BV theory of $\nIg^{(H)}$.  The positive-root height grading makes every even adjoint endomorphism strictly height-increasing and every odd adjoint endomorphism odd; hence the modular supercharacter vanishes and the reduced theory satisfies the quantum master equation.

Choose the standard complementary Lagrangian boundary polarizations, fixing the $\nIg^{(H)}[1]$ coordinate at one endpoint and its cotangent coordinate at the other.  The algebraic boundary quantization, the interval amplitude between the two augmentation states, and the holomorphic boundary current envelope are then
\[
 U(\nIg^{(H)}),\qquad
 \mathbbm1\tensor^{\mathbb L}_{U(\nIg^{(H)})}\mathbbm1
 \simeq C_*(\nIg^{(H)}),\qquad V_H.
\]
The height quotient maps preserve the polarizations and yield the inverse systems
\[
 \widehat U_h(\nIg),\qquad
 C_{*,h}^{\wedge}(\nIg),\qquad
 \widehat V_{\Delta_5}.
\]
\end{theorem}
\begin{proof}
The classical master equation is the Jacobi identity paired with invariance of the cotangent pairing.  Elliptic descent gives the classical factorization algebra.  Harmonic reduction along the interval gives a finite-dimensional BV space.  If $x$ is even of positive root height, $\operatorname{ad}_x$ raises height and is nilpotent on the finite quotient, so its restrictions to the even and odd summands both have trace zero; if $x$ is odd, $\operatorname{ad}_x$ is odd and has supertrace zero.  Thus the modular character vanishes and the finite BV quantum master equation follows.

The chosen boundary polarization quantizes the linear Kirillov--Kostant boundary bracket to $U(\nIg^{(H)})$.  Gluing the two augmentation states gives the two-sided derived bar, and the PBW antisymmetrization identifies it with $C_*(\nIg^{(H)})$.  Holomorphic descendants generate $U^{\mathrm{ch}}(\operatorname{Cur}_E\nIg^{(H)})=V_H$.  All maps preserve root height, so they pass coefficientwise to the separated inverse limits.
\end{proof}

\part{The compact scalar character}

\chapter{Reduced curve counting on \texorpdfstring{$K3\times E$}{K3 x E}}
Let $S$ be a nonsingular projective K3 surface and $E$ an elliptic curve. For a primitive class $\beta_h\in H_2(S,\Z)$ with $\beta_h^2=2h-2$, let $N_{g,h,d}$ be the reduced disconnected Gromov--Witten invariant of class $(\beta_h,d)$ with the standard incidence insertion. Define
\[
 \mathcal Z(u,q,s)=
 \sum_{g,h,d}N_{g,h,d}u^{2g-2}q^{h-1}s^{d-1}.
\]

\begin{theorem}[Primitive Igusa cusp-form theorem]
Under $r=e^{iu}$, the partition function is the Laurent expansion
\[
 \mathcal Z(u,q,s)=-\frac1{\ChiTen(r,q,s)}.
\]
Equivalently, in theta normalization,
\[
 \mathcal Z=-2^{12}\Dtheta^{-2}.
\]
\end{theorem}
\begin{proof}
Oberdieck and Pixton identify the primitive reduced partition function with the Laurent expansion of $-\ChiTen^{-1}$ \cite{OP}; their normalization is the monic Borcherds product $\ChiTen$. The second formula is the identity $\ChiTen=2^{-12}\Dtheta^2$.
\end{proof}

The theorem identifies the numerical generating series in its stated
primitive sector. The Borcherds root algebra supplies a separate
algebraic character to compare with that series.

\begin{corollary}[Exact character comparison]
Let $\jmath:\HHtwo\to\HHtwo$ be the half-coordinate map $\jmath(Z)=Z/2$, and let
\[
 \Fock(\nIg)=
 \Sym(\nIg_{\bar0})\tensor\bigwedge(\nIg_{\bar1}).
\]
Then
\[
 \mathcal Z^{\red}_{S\times E}
 =-\jmath^*\!\left(
 e^{2\rho}\sch\bigl(\Fock(\nIg)^{\tensor2}\bigr)
 \right)
 =-\ChiTen^{-1}.
\]
\end{corollary}
\begin{proof}
Put
\[
 P=\prod_{\alpha>0}(1-e^{-\alpha})^{\sdim\gIg_\alpha}.
\]
The denominator theorem gives $\den(\gIg)=e^{-\rho}P=\Dmon(2Z)$, while PBW gives $\sch\Fock(\nIg)=P^{-1}$.  Hence
\[
 e^{2\rho}\sch\bigl(\Fock(\nIg)^{\tensor2}\bigr)
 =\den(\gIg)^{-2}.
\]
Pullback by $\jmath$ sends $\den(\gIg)$ to $\Dmon(Z)$, whose square is $\ChiTen$.  The reduced incidence convention supplies the overall minus sign in the Oberdieck--Pixton theorem.
\end{proof}

The equality compares the completed root-Fock supercharacter with the
numerical reduced generating series. Its proof uses the denominator
identity and the half-coordinate pullback. A map from compact geometric
states to this Fock space, preserving operators and their trace,
is an additional comparison.

\chapter{Extensions over the elliptic classifying stack}
Work on locally noetherian complex schemes with the fppf topology.
Let \(BE\) be the stack of right \(E\)-torsors. For a left
\(E\)-stack \(M\), an object of \([M/E](T)\) is a right
torsor \(P\to T\) and an equivariant map \(P\to M\).
Forgetting the map gives \([M/E]\to BE\). All products of
stacks below are two-fibre products.

\begin{proposition}[The common translation torsor]
\label{igusa:platonic-common-torsor}
There is a natural equivalence
\[
 [(M\times N)/E_{\mathrm{diag}}]
 \simeq[M/E]\times_{BE}[N/E].
\]
The equivalences for any finite number of factors commute with
regrouping and insertion of the terminal \(E\)-stack.
\end{proposition}
\begin{proof}
An object on the right is \((P,m),(Q,n)\) together with a
torsor isomorphism \(\alpha:P\to Q\). Send it to
\((P,m,n\alpha)\). Conversely, use a common torsor twice
and its identity isomorphism. A pair of torsor maps \(a,b\)
on the right satisfies \(b\alpha=\alpha'a\), so \(b\)
is determined by \(a\). The same construction identifies the
maps and their descent isomorphisms. It commutes with base change.
For several factors, every regrouping gives one torsor with the
same list of equivariant maps, proving all stated coherences.
\end{proof}

For \(M=N=E\) with translation, the diagonal quotient is
\(E\), with coordinate \(y-x\). Both separate quotients
are points. Their fibre product over \(BE\) retains the
relative translation, while their ordinary product loses it.

Let \(\mathfrak M\) be the stack of coherent sheaves on
\(X=S\times E\), with families flat over the parameter
scheme, and put \(\mathfrak Q=[\mathfrak M/E]\).
An object of \(\mathfrak Q(T)\) is a torsor \(P\to T\)
and a flat sheaf on \(S\times P\). Fix a polarization and
write \(\mathfrak Q_\alpha\) for a Hilbert-polynomial component.
Translation preserves these polynomials. Let
\(\mathfrak H_{\alpha,\beta}\) classify short exact sequences
\(0\to A\to F\to B\to0\) on a common \(S\times P\),
with all terms flat over \(T\) and polynomials
\(\alpha,\alpha+\beta,\beta\). Its span over \(BE\) is
\[
 \mathfrak Q_\alpha\times_{BE}\mathfrak Q_\beta
 \xleftarrow{p}\mathfrak H_{\alpha,\beta}
 \xrightarrow{q}\mathfrak Q_{\alpha+\beta}.
\]
The category of spans of fppf stacks over \(BE\) has composition
by fibre product and tensor product \(\times_{BE}\). Its unit
object is the identity map of \(BE\).

\begin{theorem}[Geometric Hall correspondence]
\label{igusa:platonic-common-torsor-hall}
For fixed \(\alpha,\beta\), \(q\) is representable by
algebraic spaces and proper. The spans form a coherently associative
unital Hall object in the category of spans over \(BE\).
Its unit is the zero-sheaf span
\(BE\leftarrow BE\to\mathfrak Q\).
\end{theorem}
\begin{proof}
Over a fixed \((P,F)\), the fibre of \(q\) classifies
quotients \(F\twoheadrightarrow B\) of polynomial \(\beta\).
Its kernel is flat, because \(F\) and \(B\) are flat.
An automorphism fixing the quotient and \(F\) is the identity.
Trivializing \(P\) fppf-locally identifies this fibre with the
fixed-polynomial Quot scheme, projective over the parameter scheme
by Grothendieck's theorem \cite[Theorem~3.2]{IgusaGrothendieckQuot1961}.
The quotient functor descends to an algebraic space, and properness
descends with it.

Both composites of three Hall spans classify
\(0\subset A\subset D\subset F\) on the same torsor.
One bracketing concatenates inclusions. The other starts with
\(A\subset F\) and a subobject \(B\subset F/A\), and sets
\(D=F\times_{F/A}B\). These constructions are inverse on
objects and isomorphisms. They preserve flatness and commute with
base change. For longer flags all regroupings identify the same
subobjects and quotients. Their comparison maps agree, proving the
pentagon and higher coherences. Finally \(\mathfrak Q_0=BE\).
The sequences with zero subobject or zero quotient give the unit
and its coherences.
\end{proof}

Properness here concerns the full coherent-sheaf stack on a fixed
pair of Hilbert-polynomial components. A stability restriction can
replace its fibre by an open part of a Quot scheme, which need not
be proper. A bounded charge window also needs closure under the
extensions and flags used in its operations. A reduced critical-coefficient realization must supply pullback,
pushforward, and orientation comparisons on these same flags. Its unit has
stabilizer \(E\). A coefficient theory for this quotient must
include that nonaffine stabilizer, or work relative to \(BE\),
on the unit and every flag.


\chapter{Additive physical charge and quadratic automorphic degree}
Let
\[
 \Lam=\widetilde H(S,\Z)
\]
be the Mukai lattice of signature $(4,20)$. The electric--magnetic charge lattice is
\[
 \Gammphys=\Lam\oplus\Lam,
 \qquad c=(Q,P).
\]
Its addition is the physical superselection law. Define the Gram map
\[
 \PiGram(Q,P)=
 \left(\frac{Q^2}{2},Q\cdot P,\frac{P^2}{2}\right)
 \in\Gammaind.
\]

\begin{proposition}[Polarization identity]
For $c_i=(Q_i,P_i)$,
\[
 \PiGram(c_1+c_2)=\PiGram(c_1)+\PiGram(c_2)+B(c_1,c_2),
\]
where
\[
 B(c_1,c_2)=
 \bigl(Q_1\cdot Q_2,
 Q_1\cdot P_2+Q_2\cdot P_1,
 P_1\cdot P_2\bigr).
\]
The cross-term is a symmetric vector-valued two-cocycle:
\[
 B(c_1,c_2)+B(c_1+c_2,c_3)
 =B(c_2,c_3)+B(c_1,c_2+c_3).
\]
\end{proposition}
\begin{proof}
Expand the three quadratic expressions. The cocycle identity is bilinearity of $B$.
\end{proof}

\section{The polarized central extension}
Define the abelian central extension
\[
 0\longrightarrow\Gammaind\longrightarrow
 \widetilde\Gamma_\Pi
 \longrightarrow\Gammphys\longrightarrow0,
 \qquad
 \widetilde\Gamma_\Pi=\Gammphys\times\Gammaind,
\]
with multiplication
\[
 (c,\lambda)\cdot(d,\mu)
 =\bigl(c+d,\lambda+\mu+B(c,d)\bigr).
\]
The cocycle identity makes this multiplication associative. The polarization identity makes the graph map
\[
 j_\Pi:\Gammphys\longrightarrow\widetilde\Gamma_\Pi,
 \qquad c\longmapsto(c,\PiGram(c))
\]
a homomorphism.

Let $A=\bigoplus_{c\in\Gammphys}A_c$ be a charge-graded associative algebra, and let $\kfield[\Gammaind]$ have monomials $u^\lambda$. Define the polarized central-extension algebra
\[
 \mathcal R_\Pi(A)=A\tensor\kfield[\Gammaind]
\]
with product
\[
 (a_cu^\lambda)\star(b_du^\mu)
 =a_cb_d\,u^{\lambda+\mu+B(c,d)}.
\]
Give $a_cu^\lambda$ degree $(c,\lambda)\in\widetilde\Gamma_\Pi$.

\begin{theorem}[Graph-multiplicative polarization]
The product $\star$ is associative and $\widetilde\Gamma_\Pi$-graded. The map
\[
 \iota_\Pi:A\longrightarrow\mathcal R_\Pi(A),
 \qquad a_c\longmapsto a_cu^{\PiGram(c)}
\]
is an algebra homomorphism whose homogeneous degree is the graph element $j_\Pi(c)$.
\end{theorem}
\begin{proof}
The two parenthesizations differ by
\[
 B(c,d)+B(c+d,e)-B(d,e)-B(c,d+e),
\]
which vanishes by the cocycle identity. The polarization formula gives
\[
 \iota_\Pi(a_c)\star\iota_\Pi(b_d)
 =a_cb_d\,u^{\PiGram(c+d)}.
\]
Thus the graph of the quadratic Gram map is multiplicative inside the central extension.
\end{proof}

The same construction applies to Lie superalgebras and their enveloping algebras. It retains additive physical charge, Gram degree, and the interaction cross-term in one grading group.

\chapter{Point-local and wrapped sectors}
Let $p_E:S\times E\to E$ be projection.

\begin{lemma}[Positive elliptic degree is wrapped]
Let $C\subset S\times E$ be a proper one-dimensional subscheme whose pushforward $(p_E)_*[C]$ has positive degree. Then $p_E(C)=E$.
\end{lemma}
\begin{proof}
Properness makes $p_E(C)$ closed. Positive pushforward degree implies that at least one irreducible component of $C$ dominates a one-dimensional closed subset of $E$. Every one-dimensional closed subset of the irreducible curve $E$ is $E$ itself, so the image of that component, and hence the image of $C$, is all of $E$.
\end{proof}

Projection-finite objects define point-supported sectors on $\Ran(E)$. Positive elliptic degree defines wrapped sectors supported over the entire elliptic curve.

\begin{definition}[The local--wrapped factorization operad]
Let $\mathsf{Ch}_E$ denote the chosen chiral/factorization operad on $E$ and let $\mathsf{Ass}$ denote the ordered associative operad. Define the colored operad $\mathsf{Hyb}_E$ with colors
\[
 \mathsf l,\qquad \mathsf w_m\quad(m\ge1)
\]
by the following operation objects.
\begin{enumerate}
\item Operations with local output are
\[
 \mathsf{Hyb}_E(\mathsf l^{\otimes k};\mathsf l)=\mathsf{Ch}_E(k).
\]
\item For $r\ge1$ and $m=m_1+\cdots+m_r$, operations with wrapped output are
\[
 \mathsf{Hyb}_E
 (\mathsf l^{\otimes k},\mathsf w_{m_1},\ldots,\mathsf w_{m_r};\mathsf w_m)
 =\mathsf{Ch}_E(k)\tensor\mathsf{Ass}(r).
\]
All other operation objects are initial.
\end{enumerate}
Operadic composition uses chiral insertion and disjoint union in the first factor and ordered substitution in the second. The equality
\[
 (\gamma_1\circ\gamma_2)\tensor(a_1\circ a_2)
 = (\gamma_1\tensor a_1)\circ(\gamma_2\tensor a_2)
\]
is the mixed interchange law. The arities $(k,r)=(k,0),(0,2),(k,1)$ recover, respectively, local factorization, wrapped multiplication, and the action of local insertions on a wrapped sector.
\end{definition}

\begin{definition}[Hybrid local--wrapped algebra]
A hybrid local--wrapped algebra over $E$ is an algebra over $\mathsf{Hyb}_E$. It consists of a factorization algebra $\F_{\loc}$, a height- or charge-complete family $W=\prod_{m>0}W_m$, ordered products
\[
 \mu_{m_1,\ldots,m_r}:W_{m_1}\tensor\cdots\tensor W_{m_r}
 \longrightarrow W_{m_1+\cdots+m_r},
\]
and factorization-compatible actions of the local operation objects $\mathsf{Ch}_E(k)$ on every wrapped product. The operad axioms state associativity, local factorization, and compatibility of the two operations.
\end{definition}

A compact Hall theory of $S\times E$ produces a correspondence-valued $\mathsf{Hyb}_E$-algebra when disjoint-support direct sums, wrapped extensions, and the mixed extension correspondences satisfy the same base-change identities.


\part{The Borcherds algebra and its quantum group}

\chapter{The local--wrapped root algebra}
The positive cone is graded by $\gamma=(n,\ell,m)$, with elliptic degree $m$.  Set
\[
 \nIg_{\loc}=\bigoplus_{\substack{\alpha>0\\m(\alpha)=0}}\gIg_\alpha,
 \qquad
 \nIg_{\wrap}=\bigoplus_{\substack{\alpha>0\\m(\alpha)>0}}\gIg_\alpha.
\]

\begin{proposition}[Semidirect local--wrapped decomposition]
The space $\nIg_{\loc}$ is a Lie subsuperalgebra, $\nIg_{\wrap}$ is a Lie ideal, and
\[
 \nIg=\nIg_{\loc}\ltimes\nIg_{\wrap}.
\]
The proper height filtration preserves both summands.
\end{proposition}
\begin{proof}
The bracket adds root degrees.  The sum of two roots of elliptic degree zero again has degree zero, and addition of a nonnegative degree to a positive degree remains positive.  Height is additive on every nonzero bracket.
\end{proof}

Define the complete enveloping Hopf algebra and current envelope
\[
 \mathbb U_{\Delta_5}=\widehat U_h(\nIg),
 \qquad
 \mathbb V_{\Delta_5,E}:=\widehat V_{\Delta_5}=\varprojlim_HU^{\ch}\!\left(\Cur_E(\nIg^{(H)})\right).
\]
The adjoint action of $\nIg_{\loc}$ on $\nIg_{\wrap}$ extends by derivations to $\widehat U_h(\nIg_{\wrap})$.  Let $\# $ denote the corresponding completed smash product.

\begin{definition}[Current-generated local--wrapped operad]
The colored operad $\mathsf{Hyb}^{\cur}_{E,\Delta}$ is generated by the chiral current operations of $\Cur_E(\nIg_{\loc})$, the ordered multiplications of $\widehat U_h(\nIg_{\wrap})$, and the adjoint derivation action of the former on the latter, modulo the chiral Jacobi, associative, derivation, and semidirect-product relations.
\end{definition}

\begin{theorem}[Local--wrapped Borcherds decomposition]
Let $\nIg_{\loc}$ be the sum of positive root spaces of elliptic degree $m=0$ and let $\nIg_{\wrap}$ be the sum of those with $m>0$.  Then
\[
 \nIg=\nIg_{\loc}\ltimes\nIg_{\wrap},
\]
where $\nIg_{\loc}$ is a Lie subsuperalgebra and $\nIg_{\wrap}$ is a Lie ideal.  For the proper-height completion,
\[
 \widehat U_h(\nIg)
 \cong
 \widehat U_h(\nIg_{\wrap})\#\widehat U_h(\nIg_{\loc}),
\]
and
\[
 \widehat C_{*,h}(\nIg)
 \simeq
 \Tot C_*(\nIg_{\loc};\widehat C_{*,h}(\nIg_{\wrap})).
\]
The primitive Lie superalgebra of the completed enveloping algebra is the height completion $\widehat{\nIg}_h=\varprojlim_H\nIg^{(H)}$, and the Weyl-shifted Euler supercharacter of the completed Chevalley complex is $\Dmon(2Z)$.
\end{theorem}
\begin{proof}
Root addition adds elliptic degree.  Thus the $m=0$ summand is closed under the bracket, and every bracket with an $m>0$ root again has $m>0$.  PBW orders wrapped generators before local generators and gives the completed smash product.  The Chevalley complex of a semidirect product is the Hochschild--Serre total complex.  Primitive elements of a completed enveloping algebra are recovered height by height.  The denominator theorem supplies the last character identity.
\end{proof}

\begin{remark}
The theorem constructs a Borcherds root algebra with local and wrapped colors.  The word Hall is reserved for an algebra obtained from geometric extension correspondences.
\end{remark}

\chapter{The formal Igusa quantum group}
Let $\mathfrak b_+$ be the positive Borcherds Borel half.  Form the restricted graded dual $\mathfrak b_-:=\mathfrak b_+^{\vee,\mathrm{gr}}$ with bracket and cobracket transported by the invariant pairing, identify the two Cartan copies along their common radical quotient, and complete by a proper root height.  The completed Drinfeld double of this paired Borel datum is denoted $\gIg^{\mathrm{pair}}$.  Its positive and negative halves, diagonal Cartan, invariant form, and canonical tensor form a height-admissible Manin triple and hence a quasitriangular Lie superbialgebra.

\begin{definition}[Height-admissible Lie superbialgebra]
A root-graded Lie superbialgebra is height-admissible when every root space is finite-dimensional, every bounded-height region contains finitely many roots, and the bracket, cobracket, invariant tensor, and all iterated operations are coefficientwise finite in the height completion.
\end{definition}

\begin{proposition}[Height admissibility]
The paired Gritsenko--Nikulin Borcherds Lie superbialgebra $\gIg^{\mathrm{pair}}$ is height-admissible.
\end{proposition}
\begin{proof}
The Borcherds root spaces are finite-dimensional and the chosen height is proper.  The bracket and cobracket preserve total root degree.  In a bounded total height, finitely many decompositions of a root occur.  The canonical tensor is therefore a coefficientwise finite sum in every height quotient.
\end{proof}

\begin{theorem}[Quantized Borcherds superalgebra]
Choose the symmetrizable Borcherds--Cartan superdatum and multiplicity index set that present $\gIg$.  Hong's construction gives a quantized enveloping Hopf superalgebra
\[
 U_q(\gIg)
\]
over $\Q(q)$, generated by its Cartan elements and the real and imaginary simple-root generators, with triangular decomposition
\[
 U_q(\gIg)\cong U_q^-(\gIg)\tensor U_q^0(\gIg)\tensor U_q^+(\gIg).
\]
Its PBW root spaces have the same graded ranks as the classical Borcherds superalgebra.  The canonical pairing of the two Borel halves yields a universal $R$-matrix in the root-height completion of
\[
 U_q^+(\gIg)\widehattensor U_q^-(\gIg).
\]
Consequently the quantization preserves the root-multiplicity data that enter $\Dmon(2Z)$.
\end{theorem}
\begin{proof}
The Gritsenko--Nikulin algebra is a symmetrizable Borcherds--Kac--Moody Lie superalgebra.  Apply the generators-and-relations construction of the quantized Borcherds superalgebra.  The nondegenerate bilinear form, triangular decomposition, centre, and universal $R$-matrix are established by Hong \cite{Hong98}.  The PBW theorem identifies the associated graded root spaces with their classical counterparts.  Since the denominator exponents are the classical signed root multiplicities, they are unchanged by the flat root-graded deformation.
\end{proof}

\begin{remark}
The quantum group is attached to the Borcherds generalized Cartan superdatum.  A geometric Hall quantum group is identified with it through a Hall coproduct, pairing, completion, and primitive comparison.
\end{remark}

\chapter{A holomorphic--topological BF realization}
For every height bound $H$, let $\nIg^{(H)}$ be the finite-dimensional nilpotent quotient.  On $E\times\R$ form the shifted cotangent theory of
\[
 \Lcal_H=\Omega^{0,*}(E)\widehattensor\Omega^*(\R)\tensor\nIg^{(H)}
\]
with action
\[
 S_H(A,B)=\int_{E\times\R}\Omega_E\wedge
 \left\langle B,(\dbar+\dd_{\R})A+\frac12[A,A]\right\rangle .
\]

\begin{theorem}[Classical local and exact zero-mode realization]
For every $H$:
\begin{enumerate}
\item $S_H$ satisfies the classical master equation and defines a classical holomorphic--topological factorization algebra;
\item harmonic reduction along a compact interval gives the finite-dimensional cotangent BV theory of $\nIg^{(H)}$, whose positive-root grading makes its modular supercharacter zero and hence gives the quantum master equation;
\item after choosing the complementary Lagrangian boundary polarizations, the boundary algebra is $U(\nIg^{(H)})$, the interval amplitude is $\mathbbm1\tensor^{\mathbb L}_{U(\nIg^{(H)})}\mathbbm1\simeq C_*(\nIg^{(H)})$, and the holomorphic boundary current envelope is $V_H$.
\end{enumerate}
The transition maps yield a pro-object with boundary algebra $\widehat U_h(\nIg)$, interval complex $C_{*,h}^{\wedge}(\nIg)$, and current envelope $\widehat V_{\Delta_5}$.
\end{theorem}
\begin{proof}
The classical master equation is the Jacobi identity and invariance of the cotangent pairing.  The classical observables satisfy factorization descent.  Harmonic reduction is finite-dimensional.  An even adjoint operator raises positive-root height and is nilpotent on the finite quotient; an odd adjoint operator is odd.  Both have supertrace zero, so the modular term in the finite BV Laplacian vanishes.  The chosen polarization quantizes the boundary Lie--Poisson bracket to the enveloping algebra; the two augmentation states give the two-sided bar, identified with Chevalley chains by PBW.  Holomorphic descendants give the current algebra.  All structure maps preserve height.
\end{proof}

The theorem gives the precise field-theoretic realization used in this volume: a classical local theory together with an exact quantum protected zero-mode sector.

\chapter{The determinant and Fredholm square-root operators}
Let
\[
 \Vcal_H=\bigoplus_{\substack{\alpha>0\\h(\alpha)\le H}}\gIg_\alpha
\]
and define
\[
 \mathscr D_H|_{\gIg_\alpha}=(1-e^{-\alpha})\id.
\]

\begin{theorem}[Universal determinant section]
The finite-height Berezinian is
\[
 \Ber(\mathscr D_H)
 =\prod_{\substack{\alpha>0\\h(\alpha)\le H}}
 (1-e^{-\alpha})^{\sdim\gIg_\alpha}.
\]
The coefficientwise limit satisfies
\[
 e^{-\rho}\Ber(\mathscr D_\infty)=\Dmon(2Z).
\]
After pullback by $\jmath(Z)=Z/2$,
\[
 \jmath^*\!\left(e^{-\rho}\Ber(\mathscr D_\infty)\right)=\Dmon(Z),
 \qquad \Dmon(Z)^2=\ChiTen(Z).
\]
\end{theorem}
\begin{proof}
The Berezinian of a diagonal endomorphism is the product of its eigenvalues with exponent equal to the superdimension of the eigenspace.  Proper height makes every coefficient stationary.  The denominator theorem supplies the Weyl line and the root-to-Siegel coordinate map.
\end{proof}

\section{Trace-class realization on the root Hilbert superspace}
Choose homogeneous orthonormal bases of the finite-dimensional root spaces and set
\[
 \Hcal_{\Delta}=\widehat\bigoplus_{\alpha>0}\gIg_\alpha
\]
with its induced $\Z/2$ grading.  Let $\mathfrak C_{\mathrm{abs}}$ be the subchamber on which
\[
 \sum_{\alpha>0}\dim\gIg_\alpha\,
       \abs{e^{-\alpha(Z)}}<\infty.
\]
For $Z\in\mathfrak C_{\mathrm{abs}}$, define the trace-class diagonal operator
\[
 K_Z|_{\gIg_\alpha}=e^{-\alpha(Z)}\id,
 \qquad D_Z=1-K_Z.
\]
On each parity summand form the skew doubling
\[
 \mathbb A_{Z,\bar\epsilon}=
 \begin{pmatrix}0&D_{Z,\bar\epsilon}\\-D_{Z,\bar\epsilon}^{\mathsf t}&0\end{pmatrix}.
\]
Orient every finite-height truncation by the fixed root order and normalize its Pfaffian to $1$ at $K_Z=0$.  The oriented Fredholm super-Pfaffian is the stabilized ratio
\[
 \operatorname{sPf}_{F}(\mathbb A_Z)
 =\frac{\Pf_F(\mathbb A_{Z,\bar0})}
        {\Pf_F(\mathbb A_{Z,\bar1})}.
\]

\begin{theorem}[Polarized Fredholm realization]
Choose the protected parity refinement and the root-space Hilbert norms above.  For every $Z$ in the absolute-convergence chamber, the diagonal operator $K_Z|_{\gIg_\alpha}=e^{-\alpha(Z)}\id$ is trace class on the even and odd root Hilbert spaces.  For $D_Z=1-K_Z$,
\[
 \sdet_F(D_Z)
 =\frac{\det_F(D_Z|_{\bar0})}{\det_F(D_Z|_{\bar1})}
 =\prod_{\alpha>0}(1-e^{-\alpha(Z)})^{\sdim\gIg_\alpha}.
\]
After choosing the standard polarization of $\Hcal_\Delta\oplus\Hcal_\Delta^*$ and its orientation, the paritywise skew block operators $\mathbb A_{Z,\bar0}$ and $\mathbb A_{Z,\bar1}$ have an oriented Fredholm super-Pfaffian
\[
 \operatorname{sPf}_F(\mathbb A_Z)=\frac{\Pf_F(\mathbb A_{Z,\bar0})}{\Pf_F(\mathbb A_{Z,\bar1})}
 =\sdet_F(D_Z).
\]  Hence
\[
 e^{-\rho}\sdet_F(D_Z)=\Dmon(2Z).
\]
\end{theorem}
\begin{proof}
The defining summability condition for \(\mathfrak C_{\mathrm{abs}}\)
is precisely the trace norm of the diagonal operator, summed in both
parities.  The Fredholm determinant of $1-K_Z$ is the convergent product over eigenvalues; taking the even/odd ratio gives the superdeterminant.  On each finite root-height truncation the Pfaffian of the oriented skew doubling equals the determinant of $D_Z$.  Trace-norm convergence of the truncations passes this identity to the Fredholm limit.  The final equality is the Gritsenko--Nikulin denominator formula.
\end{proof}

\begin{remark}
The polarization and orientation in this theorem belong to the root Hilbert space.  A compact moduli-space orientation is a separate line bundle and is compared with this one by a specified intertwiner.
\end{remark}

This constructs the square-root operator on the root Hilbert superspace
where its stated summability condition holds. The nonemptiness estimate
below concerns the minimal virtual superspace, whose dimensions are
\(\lvert\sdim\gIg_\alpha\rvert\). Applying that estimate to the
full root superspace requires a bound on both parity dimensions.
A compact geometric realization requires an operator intertwiner and
a compatible map of determinant lines.

\chapter{What the automorphic character determines}
\begin{theorem}[Denominator data and reconstruction data]
The Borcherds product determines:
\begin{enumerate}
\item the signed root multiplicity $\sdim\gIg_\alpha$ in every active root degree;
\item the Weyl vector and denominator character;
\item the divisor and automorphic multiplier of the section.
\end{enumerate}
A Lie bracket is supplied by the Borcherds--Kac--Moody presentation, parity dimensions by a parity-refined root datum, Hall multiplication by extension correspondences, and an oriented square root by a square-root line with monodromy.  These reconstruction data refine the scalar product independently.
\end{theorem}
\begin{proof}
The first three items are read from the exponent, prefactor, and automorphy law.  Independence is witnessed by explicit pairs: two brackets on the same graded vector space have the same PBW character; an even--odd pair has signed dimension zero; distinct Hall correspondences may have the same Euler product; and twisting a square-root line by a sign local system preserves its square.  Hence each refined object is reconstructed from the additional datum named in the statement.
\end{proof}

\begin{recognition}[Complete primitive recognition]
Let $\mathfrak p$ be a positive graded Lie superalgebra and let $\Psi:\nIg\twoheadrightarrow\mathfrak p$ be a graded Lie map.  The map is an isomorphism if either:
\begin{enumerate}
\item the parity-refined Hilbert series of $U(\nIg)$ and $U(\mathfrak p)$ agree; or
\item each primitive degree has prescribed pure parity and the signed PBW characters agree.
\end{enumerate}
The induced height completions, Chevalley complexes, and current envelopes are then isomorphic.
\end{recognition}
\begin{proof}
Choose the first height containing a homogeneous kernel vector.  All decomposable PBW terms in that height are determined by lower heights and agree.  The primitive parity dimension therefore drops, contradicting the first hypothesis.  In the second case the logarithm of the PBW product recovers signed primitive dimensions inductively, and prescribed parity recovers ordinary dimensions.  Naturality of completion, Chevalley chains, and current envelopes gives the final assertion.
\end{proof}

\part{The universal Borcherds Hall and three-Calabi--Yau models}

\chapter{Finite simple-root windows}
Let $I_\Delta$ be the multiset of simple roots of the Gritsenko--Nikulin Borcherds superalgebra.  A repeated imaginary simple root is written once for each copy.  The parity map is
\[
 p:I_\Delta\longrightarrow\Z/2.
\]
Choose a proper height function
\[
 h:I_\Delta\longrightarrow\Z_{>0}
\]
whose sublevel sets are finite.  Put
\[
 I_H=\set{i\in I_\Delta:h(i)\le H}.
\]
The restriction of the invariant form gives a finite symmetric Borcherds--Cartan matrix
\[
\begin{aligned}
 A_H&=(a_{ij})_{i,j\in I_H},\\
 a_{ii}&=2 &&\text{for real }i,\\
 a_{ii}&\leq0 &&\text{for imaginary }i,\\
 a_{ij}&\leq0 &&\text{for }i\neq j.
\end{aligned}
\]
The evenness of the root lattice makes every imaginary diagonal entry even.

\begin{construction}[The Igusa Borcherds quiver]
Order $I_H$.  For $i<j$ place $-a_{ij}$ arrows from $i$ to $j$.  At an imaginary vertex $i$ place
\[
 g_i=1-\frac{a_{ii}}2
\]
loops.  A real vertex has loop number zero.  Denote the resulting quiver by $Q_{\Delta,H}$.
\end{construction}

\begin{proposition}[Euler form]
The symmetrized Euler form of $Q_{\Delta,H}$ is $A_H$:
\[
 \angles{e_i,e_j}+\angles{e_j,e_i}=a_{ij}.
\]
\end{proposition}
\begin{proof}
For $i\ne j$, the total number of arrows between the two vertices is $-a_{ij}$, so the symmetrized Euler entry is $a_{ij}$.  At $i$ the value is $2-2g_i=a_{ii}$.
\end{proof}

The inclusions $I_H\subset I_{H+1}$ identify $Q_{\Delta,H}$ with a full subquiver of $Q_{\Delta,H+1}$.  Representations supported on the old vertices form an exact extension-closed subcategory.  This gives compatible Hall embeddings.


\chapter{The real-root seed window}
The three real simple roots are
\[
 \delta_1=2f_2-f_3,
 \qquad
 \delta_2=2f_{-2}-f_3,
 \qquad
 \delta_3=f_3,
\]
with Gram matrix
\[
 A_{\mathrm{re}}=
 \begin{pmatrix}
 2&-2&-2\\
 -2&2&-2\\
 -2&-2&2
 \end{pmatrix}.
\]
Choose the order $\delta_1<\delta_2<\delta_3$.  The real-root quiver $Q_{\mathrm{re}}$ has two arrows from vertex $i$ to vertex $j$ for each $i<j$.

\begin{proposition}[Real-root Hall seed]
The symmetrized Euler form of $Q_{\mathrm{re}}$ is $A_{\mathrm{re}}$.  The Hall composition algebra generated by its three simple representations is the positive quantum Kac--Moody algebra with Cartan matrix $A_{\mathrm{re}}$.
\end{proposition}
\begin{proof}
The diagonal Euler entries are $2$.  The total number of arrows between distinct vertices is two, hence the symmetrized off-diagonal entries are $-2$.  The Ringel--Green Hall theorem identifies the composition algebra with the positive quantum algebra defined by that symmetric Cartan matrix.
\end{proof}

The three vertices already contain the Weyl group and the reflective divisor.  Imaginary vertices add the automorphic correction.  The finite-height tower therefore begins with a completely explicit hyperbolic Hall algebra and enlarges it by one finite collection of imaginary simple roots at each stage.

\section{The first Ginzburg differential}
Let $a_{ij}^{(1)},a_{ij}^{(2)}:i\to j$ be the two arrows for $i<j$.  Add reverse arrows $(a_{ij}^{(r)})^*:j\to i$ of degree $-1$ and loops $t_i$ of degree $-2$.  The canonical three-Calabi--Yau completion has
\[
 d a_{ij}^{(r)}=0,
 \qquad
 d(a_{ij}^{(r)})^*=0,
\]
and
\begin{equation}\label{eq:vol3-ginzburg-differential}
 d t_i=
 \sum_{a:s(a)=i}aa^*
 -\sum_{a:t(a)=i}a^*a.
\end{equation}
The identity $d^2=0$ follows immediately.  Formula \eqref{eq:vol3-ginzburg-differential} is the derived moment-map equation for the doubled real-root quiver.

\begin{calculation}[Vertex one]
At vertex $1$,
\[
 d t_1=
 \sum_{r=1}^2
 \left(
 a_{12}^{(r)}(a_{12}^{(r)})^*
 +a_{13}^{(r)}(a_{13}^{(r)})^*
 \right).
\]
At vertex $3$,
\[
 d t_3=-\sum_{r=1}^2\left(
 (a_{13}^{(r)})^*a_{13}^{(r)}+
 (a_{23}^{(r)})^*a_{23}^{(r)}\right).
\]
The sum over vertices is $\sum_a(aa^*-a^*a)$.
It vanishes in the cyclic quotient and after matrix trace,
but generally not in the path algebra.
\end{calculation}
\begin{proof}
The outgoing arrows at vertex $1$ have targets $2$ and $3$.
The incoming arrows at vertex $3$ have sources $1$ and $2$.
The displayed formulas follow from these incidence lists.
Each arrow contributes one commutator to the sum over vertices.
The trace of a commutator vanishes, giving the asserted trace identity.
\end{proof}


\chapter{Hall realization of the positive algebra}
Let $\mathcal A_H=\Rep_{\mathbb F_q}(Q_{\Delta,H})$.  Its twisted Hall product is
\[
 [M]*[N]
 =v^{\angles{\dim M,\dim N}}
 \sum_{[L]}
 \abs{\set{N'\subset L:N'\cong N,\ L/N'\cong M}}
 [L],
 \qquad v^2=q.
\]

The class $[0]$ is the unit: the required subobject is unique on either
side. Associativity follows by counting flags $N_1\subset N_2\subset L$
in two orders. The twist factors agree because the Euler form is
bilinear:
\[
 \langle\alpha,\beta\rangle+\langle\alpha+\beta,\gamma\rangle
 =\langle\beta,\gamma\rangle+\langle\alpha,\beta+\gamma\rangle.
\]
Thus this normalization defines the unital twisted Hall algebra on
isomorphism classes.

\begin{theorem}[Finite-window composition algebra]
At each vertex $i$ of $Q_{\Delta,H}$, choose the one-dimensional
representation $S_i$ supported at $i$, with every loop acting by zero.
Use the generic composition algebra generated by these classes across
finite fields, with the twisted subobject-count multiplication above.
After extending coefficients from $\mathbb Z[v,v^{-1}]$ to
$\mathbb C(v)$, this is the positive quantum generalized Kac--Moody
algebra for the even Borcherds--Cartan matrix $A_H$.
The parameter in Kang--Schiffmann is $v_{\mathrm{KS}}=v^{-1}$.
Their specified semisimple perverse-sheaf construction on nilpotent
representation varieties gives the corresponding canonical basis.
\end{theorem}
\begin{proof}
This is the selected-simple composition construction of
Kang--Schiffmann, Section~2.3 and Theorem~2.2, with one selected
zero-loop simple at each vertex. Their geometric realization uses
its split graded Grothendieck group over
$\mathbb Z[v_{\mathrm{KS}},v_{\mathrm{KS}}^{-1}]$
(Theorem~4.1), with the stated inversion of parameter
\cite{KangSchiffmann}. The theorem does not use every simple object
of the full representation category. A vertex with $c$ loops has
$q^c$ one-dimensional simple representations over $\mathbb F_q$,
whereas only the chosen zero-loop representation is the generator here.
\end{proof}

A semi-derived Hall realization of a full quantum group additionally
specifies two-periodic complexes, the invertible acyclic classes and
Cartan reduction \cite{LuBorcherdsHall,AxtellLee}.
Its coproduct, pairing and completion are separate from the positive
composition algebra used here.

\section{Parity and the quantized superdatum}
The Gritsenko--Nikulin denominator algebra carries a parity map on its simple-root multiset.  It is additive on the free lattice of labelled simple generators.
Projection to the Lorentzian degree can identify contributions of
different parity. Both parity dimensions must therefore be retained
in each projected root space.  Hong's construction applied to this Borcherds superdatum gives the positive Borel
\[
 U_q^+(\gIg^{(H)})
\]
and the full quantized Borcherds superalgebra $U_q(\gIg^{(H)})$.

\begin{theorem}[Ordinary and super presentations]
The ordinary quiver Hall presentation and Hong's quantum
superalgebra use the same specified Cartan matrix and labelled
simple generators. Their relation ideals and parity conventions
are different inputs. Equality of their Cartan gradings alone does
not identify their PBW characters.
\end{theorem}
\begin{proof}
The two presentations are those of the preceding Hall construction
and Hong's superalgebra construction. The difference already occurs
for a single isotropic generator. The enveloping algebra of a
one-dimensional even abelian Lie algebra is \(\mathbb C[x]\).
Declaring its generator odd only as a graded vector space leaves
a nonzero vector \(x^2\) in degree two. The enveloping algebra
of a one-dimensional odd abelian Lie superalgebra is
\(\bigwedge(\xi)\), since its defining relation is
\(2\xi^2=[\xi,\xi]=0\). Its degree-two component is zero.
Thus parity assignment to an ordinary PBW basis does not produce
the super relation ideal or its character. A comparison must
respect these relations degree by degree.
\end{proof}

\chapter{The three-Calabi--Yau completion}
Let $kQ_{\Delta,H}$ be the path algebra. Its algebraic
three-Calabi--Yau completion is
\[
 \Gamma_{\Delta,H}=\Pi_3(kQ_{\Delta,H}).
\]
Write $J$ for the ideal generated by all arrows in its graded quiver.
Its completed version is the separate topological dg algebra
$\widehat\Gamma_{\Delta,H}=\varprojlim_n\Gamma_{\Delta,H}/J^n$.
The algebraic Calabi--Yau theorem is used for $\Gamma_{\Delta,H}$;
a topological Calabi--Yau assertion for $\widehat\Gamma_{\Delta,H}$
requires its completed bimodule category.
This is the Ginzburg model with zero potential: a degree-minus-one
reverse arrow for every arrow and a degree-minus-two loop at each
vertex. The loop differential is the signed sum of incident
arrow--reverse-arrow products.

\begin{theorem}[Borcherds--Ginzburg projections]
For each finite quiver \(Q_{\Delta,H}\), its undeformed
three-Calabi--Yau completion \(\Gamma_{\Delta,H}\) is the
Ginzburg dg algebra with zero potential. If \(H\le H'\) and
\(Q_{\Delta,H}\) is the full subquiver on a subset of vertices,
there is a surjective dg algebra map
\[
 \pi_{H'H}:\Gamma_{\Delta,H'}\longrightarrow\Gamma_{\Delta,H}.
\]
It kills the added vertices, every arrow incident to them, the
corresponding reverse arrows, and the added degree-minus-two loops.
These maps compose strictly and extend continuously to the completed algebras.
\end{theorem}
\begin{proof}
The finite Calabi--Yau assertion is the Ginzburg model of Keller's
completion theorem \cite{KellerCY,GinzburgCY}. We concatenate
paths from left to right, so \(aa^*\) is based at the source
of \(a\). Original and reverse arrows have zero differential.
For the degree-minus-two loop \(t_i\), the differential is
\[
 d t_i=\sum_{s(a)=i}aa^*-\sum_{t(a)=i}a^*a.
\]
The projection deletes exactly the terms involving an added vertex,
leaving the same sum for the full subquiver. It therefore commutes
with the differential. It preserves path length and extends to the
completed path algebra. Successive projections kill the union of
their omitted vertices and arrows, proving strict composition.
\end{proof}

The evident inclusion of graded quivers does not give the reverse
dg map. A single old vertex has loop differential zero. Adding
an arrow \(a:i\to j\) changes that differential to \(aa^*\)
at \(i\). The map fixing the old loop therefore fails to commute
with the differential.

Write \(\Gamma_\Delta=(\Gamma_{\Delta,H},\pi_{H'H})_H\)
for this pro-system of dg algebras. It is separate from the
ordinary Hall embeddings obtained by extension by zero on
representations of a full subquiver. A direct system of
Calabi--Yau dg categories compatible with those Hall embeddings
requires additional dg functors. It does not follow from the
inclusion of the displayed Ginzburg generators.

The quantum group retains its separately specified root-height
completion. Its transition maps must preserve the Cartan generators,
coproduct, and pairing. Quotienting the positive nilpotent Lie
algebra by height and truncating a full quantum group are different
operations.

\chapter{Hall states and protected superstates}
The representation varieties of $Q_{\Delta,H}$ carry the
perverse-sheaf complexes used in the composition-algebra construction.
The full equivariant derived category is a larger carrier.
Already for one real vertex and dimension one, its quotient is
$B\mathbb G_m$, with $H^*(B\mathbb G_m;\mathbb Q)=\mathbb Q[c_1]$,
where $|c_1|=2$. Its total equivariant cohomology is infinite-dimensional.
Finite root multiplicities therefore refer to the specified composition
algebra, not to total equivariant cohomology of this stack.

\begin{definition}[Igusa Hall category]
For every dimension vector $\mathbf d$, let
\[
 \mathfrak R_{H,\mathbf d}=
 [\operatorname{Rep}_{\mathbf d}(Q_{\Delta,H})/G_{\mathbf d}].
\]
The finite-height Hall category is the direct sum of the equivariant constructible derived categories on the $\mathfrak R_{H,\mathbf d}$, equipped with Hall convolution along the stack of short exact sequences.  Denote it by $\mathcal B_{\Delta,H}$ and set
\[
 \mathcal B_\Delta=\varinjlim_H\mathcal B_{\Delta,H}.
\]
\end{definition}

Independently, the protected one-particle superspace is the root-graded superspace
\[
 \mathcal H_{\mathrm{1p},H}
 =\bigoplus_{\alpha>0}\gIg^{(H)}_\alpha,
\]
with the parity of the Borcherds superalgebra.  Its multiparticle completion is the PBW/Fock superspace of $U(\nIg^{(H)})$.

\begin{theorem}[Hall and protected-state comparison]
The composition subalgebra in the finite-window Hall theorem
is the ordinary positive quantum generalized Kac--Moody algebra.
This statement does not identify the Grothendieck group of the full
category $\mathcal B_{\Delta,H}$ with that composition algebra.  The protected one-particle superspace has signed root multiplicities
\[
 \sdim\mathcal H_{\mathrm{1p},H,\alpha}
 =\sdim\gIg^{(H)}_\alpha.
\]
Its height-complete PBW character is
\[
 \sch\Fock(\nIg)
 =\prod_{\alpha>0}(1-e^{-\alpha})^{-\sdim\gIg_\alpha},
\]
and its Weyl-shifted inverse is $\Dmon(2Z)$.  The composition canonical basis and the protected root presentation use the specified simple-root indexing and Cartan grading; the parity refinement belongs to the Borcherds superdatum.
\end{theorem}
\begin{proof}
The composition-algebra assertion is the finite-window theorem. The protected assertion is PBW for the Borcherds Lie superalgebra.  Proper height passes both constructions to the limit.  The Gritsenko--Nikulin denominator theorem gives the final product identity.
\end{proof}

\chapter{Extension determinants and critical orientations}
Let \(\mathcal C\) be a three-Calabi--Yau dg category over
\(\mathbb C\), with its chosen trace. For families of objects
\(A,B\) over a base \(T\), let \(C(A,B)\) denote the relative
derived Hom complex on \(T\). Suppose these complexes are perfect,
and that the trace gives base-change-compatible duality
\[
 C(B,A)\simeq C(A,B)^\vee[-3].
\]
Put \(L(A,B)=\det C(A,B)\) and \(D(A)=L(A,A)\).
The determinant convention is
\(\det C=\bigotimes_i(\det C^i)^{(-1)^i}\), with symmetry sign
specified by \(\chi(C)\bmod2\).

\begin{theorem}[Extension determinant]
For an exact triangle \(A\to F\to B\to A[1]\), there is a
base-change-compatible isomorphism
\[
 D(F)\simeq D(A)\otimes D(B)\otimes L(A,B)^{\otimes2}.
\]
The determinant comparisons obtained by refining a flag agree, with
the stated determinant symmetry signs.
\end{theorem}
\begin{proof}
Apply derived Hom to the triangle in each variable. Determinant
additivity gives
\[
 D(F)\simeq L(B,A)\otimes D(A)\otimes D(B)\otimes L(A,B).
\]
Dualization and the odd shift each invert the determinant line.
Thus Calabi--Yau duality identifies \(L(B,A)\) with \(L(A,B)\).
Reorder by determinant symmetry to obtain the formula. For a full
flag, all iterated decompositions identify with the same ordered
tensor product of \(C(A_i,A_j)\), for its successive quotients
\(A_i\). Determinant additivity for refined filtrations proves
coherence. Each operation commutes with base change.
\end{proof}

A square-root orientation consists of lines \(o_A\) and
isomorphisms \(o_A^{\otimes2}\simeq D(A)\). Extension
compatibility requires maps
\[
 o_A\otimes o_B\otimes L(A,B)\longrightarrow o_F
\]
whose squares are the determinant comparisons. Choosing coordinates
on a smooth critical chart trivializes its ambient canonical line.
Hall convolution requires transition maps compatible with the cross
line and the zero object. On three-step flags, the two orientation
maps can differ by an element of \(\mu_2\), even when their
squares agree. That sign must be trivialized coherently.

\begin{example}[A nontrivial cross line]
Let \(X\) be a smooth projective Calabi--Yau threefold, with a
chosen volume form, and let \(x\in X\). On
\(X\times T\), where \(T=\mathbb P^1\), take
\[
 A=\mathcal O_X\boxtimes\mathcal O_T(1),\qquad
 B=\mathcal O_x\boxtimes\mathcal O_T.
\]
Then \(C(A,B)=\mathcal O_T(-1)[0]\), and relative Serre
duality gives \(C(B,A)=\mathcal O_T(1)[-3]\).
Thus \(L(A,B)=L(B,A)=\mathcal O_T(-1)\).
The self-Hom complexes of \(A\) and \(B\) have constant
determinant lines, since the factors \(\mathcal O_T(1)\)
cancel in the first self-Hom. Nevertheless
\[
 D(A\oplus B)=D(A)\otimes D(B)\otimes\mathcal O_T(-2).
\]
The orientation cross factor is \(\mathcal O_T(-1)\).
A sign change cannot replace this line of nonzero degree.
\end{example}

\begin{corollary}[Oriented critical coefficients]
Suppose the chosen critical charts have square-root orientation
maps on overlaps and extension flags, with unit and flag coherences.
Suppose their vanishing-cycle comparison maps preserve these data.
Then the local vanishing-cycle complexes descend to the specified
oriented coefficient object. Their Hall convolution is associative
whenever its pullback and pushforward maps are defined and satisfy
base change on the flag diagrams.
\end{corollary}
\begin{proof}
The overlap cocycle identifies the local coefficient complexes by
descent. On a three-step flag, orientation coherence identifies
the two Thom--Sebastiani coefficient maps. Proper or compact-support
base change, on the stipulated domain, identifies the two composites
of convolution. The zero flag gives the unit equation.
\end{proof}

For the quiver charts, the remaining construction is an equivariant
system of square roots and coefficient maps on the actual
representation and extension stacks, compatible with changes of
height. For compact \(S\times E\), the reduced complexes and
the common elliptic torsor must also be retained. A coordinate
volume on an affine chart specifies neither descent system by itself.

\chapter{The colored elliptic Ran space}
The elliptic degree is a global charge carried by the whole configuration.  Define
\[
 \Ran^{\mathrm{wr}}(E)=\coprod_{m\ge0}\Ran(E)_m
\]
with monoidal operation
\[
 (I,m)\sqcup(J,n)=(I\sqcup J,m+n)
\]
on disjoint supports.  The component $m=0$ is point-local.  The components $m>0$ are wrapped.

Let
\[
 \nIg=\bigoplus_{m\ge0}\nIg[m]
\]
be the elliptic-degree grading.  The root bracket obeys
\[
 [\nIg[m],\nIg[n]]\subset\nIg[m+n].
\]

\begin{construction}[Hybrid factorization algebra]
Form the chiral Chevalley factorization coalgebra
\[
 \mathcal F^{\mathrm{hyb}}_{\Delta,E}
 =C_*^{\mathrm{ch}}\!\left(\Cur_E(\nIg)\right).
\]
Retain the elliptic-degree grading on every Chevalley word and place its degree-$m$ summand on $\Ran(E)_m$.  Disjoint union adds the degree labels.
\end{construction}

\begin{theorem}[Local--wrapped factorization]
$\mathcal F^{\mathrm{hyb}}_{\Delta,E}$ is a graded factorization coalgebra on $\Ran^{\mathrm{wr}}(E)$.  Its degree-zero color is the current theory of $\nIg[0]$.  Every degree-$m$ color is a module for the degree-zero color.  The factorization product of degrees $m$ and $n$ lands in degree $m+n$.  Its chiral primitives are $\Cur_E(\nIg)$, and constant modes recover $\nIg$.
\end{theorem}
\begin{proof}
The chiral Chevalley construction is a factorization coalgebra for every chiral Lie algebra.  The bracket preserves elliptic-degree addition, so its differential and coproduct preserve the graded Ran placement.  The primitive theorem for the Chevalley coalgebra recovers the current Lie algebra.
\end{proof}

This construction retains the nonnegative root-degree label on a
point-supported current theory. It does not replace a sheaf projecting
onto the whole elliptic curve by a point-supported sheaf. The compact
wrapped comparison must map its common-torsor extension correspondences
to these graded current operations, preserving support and charge.
The bracket on the present carrier is the Borcherds bracket.

\chapter{Physical charge and Gram degree}
Let $\Gammphys$ be an additive physical charge lattice and let
\[
 \PiGram(c)=\left(\frac{Q^2}{2},Q\cdot P,\frac{P^2}{2}\right)
\]
be the Gram map.  Its polarization is
\[
 B(c,d)=\PiGram(c+d)-\PiGram(c)-\PiGram(d).
\]
Define the central extension
\[
 \widetilde\Gamma_\Pi=\Gammphys\times\Gammaind
\]
with multiplication
\[
 (c,\lambda)(d,\mu)
 =(c+d,\lambda+\mu+B(c,d)).
\]

\begin{theorem}[Graph homomorphism]
The map
\[
 j_\Pi:\Gammphys\longrightarrow\widetilde\Gamma_\Pi,
 \qquad
 c\longmapsto(c,\PiGram(c))
\]
is a group homomorphism.  For every $\Gammphys$-graded algebra $A$, the Rees lift
\[
 a_c\longmapsto a_cu^{\PiGram(c)}
\]
turns multiplication into an additive $\widetilde\Gamma_\Pi$-grading with cocycle $u^{-B(c,d)}$.
\end{theorem}
\begin{proof}
The defining multiplication gives
\[
 j_\Pi(c)j_\Pi(d)
 =(c+d,\PiGram(c)+\PiGram(d)+B(c,d))
 =j_\Pi(c+d).
\]
The algebra identity follows immediately.
\end{proof}

The universal Igusa category is graded directly by the Lorentzian root lattice.  A geometric $K3\times E$ theory enters through a charge functor to $\widetilde\Gamma_\Pi$.

\chapter{The universal protected field theory}
For each $H$ take cotangent BF theory with gauge Lie superalgebra $\nIg^{(H)}$ on $\R^2_{\mathrm{top}}\times E$.  Its fields are
\[
 \Omega^\bullet(\R^2)\widehattensor
 \Omega^{0,\bullet}(E)
 \widehattensor
 \left(\nIg^{(H)}[1]\oplus(\nIg^{(H)})^\vee[2]\right).
\]
The action is
\[
 S_H=\int_{\R^2\times E}
 \angles{B,(\dd+\dbar)A+\frac12[A,A]}.
\]

\begin{theorem}[Classical local theory and exact quantum harmonic sector]
For every finite height $H$ the shifted-cotangent action defines a classical mixed holomorphic--topological factorization algebra on $\R^2_{\mathrm{top}}\times E$.  Harmonic reduction on a compact product cell gives the finite-dimensional cotangent BV theory of $\nIg^{(H)}$, and this reduced theory satisfies the quantum master equation.  Complementary Lagrangian boundary polarizations give the boundary enveloping algebra $U(\nIg^{(H)})$, the two-sided bar complex $C_*(\nIg^{(H)})$, and the current envelope on $E$.  The proper-height inverse system is coefficientwise finite and defines a complete classical local theory together with an exact quantum harmonic sector.
\end{theorem}
\begin{proof}
The invariant cotangent pairing and the Jacobi identity give the classical master equation and factorization descent.  Harmonic reduction leaves a finite positive nilpotent Lie superalgebra.  Every even adjoint endomorphism raises proper height and is strictly upper triangular; every odd adjoint endomorphism has vanishing supertrace.  The reduced BV divergence therefore vanishes.  Canonical quantization in the selected polarization gives the enveloping algebra, and the augmentation boundary states give the two-sided bar, identified with Chevalley chains by PBW.  Each map preserves height, so the inverse system is defined coefficientwise.
\end{proof}


\chapter{Absolute convergence of the virtual root operator}
For each positive root $\alpha$, let $V_\alpha$ be the minimal parity realization of the signed multiplicity: it is purely even of dimension $\sdim\gIg_\alpha$ when this integer is nonnegative and purely odd of dimension $-\sdim\gIg_\alpha$ when it is negative.  Set
\[
 \Vcal_\Delta=\widehat\bigoplus_{\alpha>0}V_\alpha.
\]
This virtual root superspace is the canonical linear carrier of the denominator character; the Lie bracket remains carried by $\gIg$.

Let $h(\alpha)$ be the proper height.  Jacobi-form coefficient estimates give constants $C_0,C_1>0$ such that
\begin{equation}\label{eq:vol3-root-growth}
 \dim V_\alpha=\abs{\sdim\gIg_\alpha}\le C_0\exp(C_1h(\alpha))
\end{equation}
for every positive root.  The number of Gram triples of height $H$ is polynomial in $H$.

\begin{theorem}[Nonempty trace-class chamber]
Let $Z_0$ lie in the interior of the positive Weyl chamber.  For sufficiently large $t>0$,
\[
 \sum_{\alpha>0}\dim V_\alpha\,
 \abs{e^{-\alpha(tZ_0)}}<\infty.
\]
Hence the absolute-convergence chamber contains a conical tail in every interior direction.
\end{theorem}
\begin{proof}
Interior positivity gives a constant $c(Z_0)>0$ with
$\Re\alpha(Z_0)\ge c(Z_0)h(\alpha)$.
Combine this with \eqref{eq:vol3-root-growth}.  The contribution of roots of height $H$ is bounded by a polynomial in $H$ times
\[
 \exp\bigl((C_1-tc(Z_0))H\bigr).
\]
The series converges when $tc(Z_0)>C_1$.
\end{proof}

\begin{corollary}[Trace-class root Hamiltonian]
On the conical tail, the diagonal operator
\[
 K_Z\big|_{V_\alpha}=e^{-\alpha(Z)}\id
\]
is trace class on both parity summands.  The Fredholm product is holomorphic and equals the absolutely convergent Borcherds product.
\end{corollary}
\begin{proof}
The trace norm is the summable series of the preceding theorem, separated by parity.  Standard trace-class determinant theory gives holomorphy, and diagonalization gives the product over root eigenvalues.
\end{proof}


\chapter{The virtual root Fredholm operator}
Let
\[
 \mathcal H^{\mathrm{vir}}_\Delta
 =\widehat\bigoplus_{\alpha>0}V_\alpha
\]
with an orthonormal basis in each minimal parity realization.
This is the virtual root Hilbert space used in the preceding
absolute-convergence theorem.  In the absolute-convergence chamber define
\[
 K_Z\big|_{V_\alpha}=e^{-\alpha(Z)}\id,
 \qquad
 D_Z=1-K_Z.
\]

\begin{theorem}[Oriented Fredholm square root]
The operator $K_Z$ is trace class on the even and odd summands of the virtual root Hilbert space.  Its Fredholm Berezinian is
\[
 \sdet_F(D_Z)
 =\prod_{\alpha>0}(1-e^{-\alpha(Z)})^{\sdim\gIg_\alpha}.
\]
Choose an order of the positive roots, first by height and then lexicographically.  This choice orients every finite-height root space.  The skew doubling
\[
 \mathbb A_Z=
 \begin{pmatrix}0&D_Z\\-D_Z^{\mathsf t}&0\end{pmatrix}
\]
has an oriented Fredholm super-Pfaffian satisfying
\[
 \operatorname{sPf}_F(\mathbb A_Z)=\sdet_F(D_Z).
\]
Consequently
\[
 e^{-\rho}\operatorname{sPf}_F(\mathbb A_Z)=\Dmon(2Z).
\]
\end{theorem}
\begin{proof}
Absolute convergence of the Borcherds product is the trace-norm summability condition.  The determinant formula follows from the spectrum.  At finite height the Pfaffian of the oriented skew doubling is the determinant.  Trace-norm convergence passes the identity to the limit.  The denominator theorem supplies the Weyl shift.
\end{proof}

\section{Automorphic orientation}
The product defines a section of the automorphic line $\mathcal L^5\tensor\nu_5$.  On the absolute-convergence chamber, the oriented Fredholm Pfaffian is a trivialization of this line.  Analytic continuation uses the automorphy factor of $\Dmon$ and therefore glues the chamberwise Fredholm models with transition character $\nu_5$.

\begin{corollary}[Orientation character]
The automorphic orientation line of the root Dirac family is the sign local system $\nu_5$.  On every real simple reflection it has value $-1$.
\end{corollary}
\begin{proof}
The Fredholm product equals $D_5$ on the convergence chamber, and $D_5$ transforms with automorphy factor $\nu_5$.  This transformation law is the gluing rule for the automorphic orientation line.  A real simple root occurs with multiplicity one, so reflection reverses its oriented normal line.
\end{proof}

\chapter{The universal Igusa protected effective theory}
\begin{definition}[Universal Igusa completion]
The protected completion selected by the K3 half-index is the quintuple
\[
 \mathfrak K_{\Delta}=
 \left(
 \Gamma_\Delta,
 \mathcal B_\Delta,
 \mathcal H_{\mathrm{1p}},
 \mathcal F^{\mathrm{hyb}}_{\Delta,E},
 \mathbb A_Z
 \right).
\]
Its components are the pro-system of finite three-Calabi--Yau completions, the ordinary quiver Hall category, the Borcherds one-particle Lie superalgebra, the local--wrapped current factorization algebra, and the virtual root Dirac family.
\end{definition}

\begin{theorem}[Igusa protected completion]
The quintuple $\mathfrak K_\Delta$ satisfies:
\begin{enumerate}
\item its one-particle Lie superalgebra is $\gIg$;
\item its ordinary Hall presentation and Hong's quantum supergroup retain
the stated Cartan grading; comparison of their relation ideals and
parity-refined PBW spaces requires an additional map;
\item its chiral primitive algebra on $E$ is the local--wrapped current algebra of $\gIg$;
\item its raw multiparticle supercharacter is $P^{-1}$, where $e^{-\rho}P=\DenIg$; equivalently $e^{\rho}\sch\Fock(\nIg)=\DenIg^{-1}$;
\item the half-coordinate pullback of the doubled Weyl-normalized character is $-\ChiTen^{-1}$ in the Oberdieck--Pixton normalization;
\item its oriented first-order determinant is $\DenIg$ and its orientation monodromy is $\nu_5$.
\end{enumerate}
\end{theorem}
\begin{proof}
The first three statements are the Borcherds root, quiver Hall, Hong quantum-supergroup, Ginzburg, and colored-Ran theorems.  PBW gives the fourth.  The identity $\ChiTen=\Dmon^2$ and the established reduced scalar theorem give the fifth.  The Fredholm theorem gives the sixth.
\end{proof}

The tuple retains algebraic constructions associated with the Igusa
presentation and their scalar characters. A geometric type-IIA or
topological-string realization must construct compatible functors
from its oriented extension category to these targets. Primitive
states and extension amplitudes are inputs to those functors.
Their action on morphisms, convolution, orientations, and chiral
operations must also be specified and proved compatible.

\chapter{The primitive comparison and its extensions}
Let \(\mathfrak p\) be a positive-root-graded Lie superalgebra
with finite-dimensional homogeneous pieces and a proper positive
height. A primitive realization of the Igusa presentation is a
parity-preserving homogeneous map from its labelled simple-root
spaces to \(\mathfrak p\), satisfying the defining relations.

\begin{theorem}[Universal presentation map]
\label{igusa:primitive-presentation-map}
Let \(\mathfrak p\) be a positive-root-graded Lie superalgebra
over \(\mathbb C\), with finite-dimensional homogeneous pieces.
Suppose a homogeneous parity-preserving map from every simple-root
generating space of \(\nIg\) to \(\mathfrak p\) satisfies
the defining Borcherds--Serre relations. It extends uniquely to
a graded Lie-superalgebra map
\[
 f:\nIg\longrightarrow\mathfrak p.
\]
It induces continuous maps on height completions, enveloping
algebras, Chevalley chains, and level-zero currents. The map \(f\)
is an isomorphism if and only if it is surjective and the root-graded
PBW series agree in each parity.
\end{theorem}
\begin{proof}
Extend the chosen maps to the free Lie superalgebra on the
simple-root spaces. The relation hypothesis makes the defining
ideal lie in its kernel, giving \(f\). The generators prove
uniqueness. On enveloping algebras the map sends
\(x_1\cdots x_j\) to \(f(x_1)\cdots f(x_j)\).
On Chevalley chains it is the graded symmetric map on the shifted
generators. It commutes with the Chevalley differential because
\(f\) preserves brackets. The current map is
\(\partial^j x\mapsto\partial^j f(x)\), preserving
\([x_\lambda y]=[x,y]\). Root degree preservation gives
compatible maps at each height and hence on the stated limits.

If \(f\) is surjective, \(U(f)\) is surjective. Proper
positive height makes each root-graded PBW component
finite-dimensional. Equality of its even and odd dimensions makes
\(U(f)\) bijective in every component. PBW embeds each Lie
superalgebra into its enveloping algebra, so \(f\) is injective.
The converse follows by applying the constructions to its inverse.
\end{proof}


A geometric comparison also requires functors on the actual Hall
and critical categories, a compatible coproduct and Hopf pairing,
and maps of determinant lines and operator domains. These categories
and lines are not the domain of the primitive Lie map. The
additional maps must be constructed and checked against that map.
They are not consequences of its universal property.

\chapter{The eight diagonal-divisor forms}
Cl\'ery and Gritsenko classify eight genus-two modular forms whose divisor is the diagonal with order one \cite{GC}.  Gritsenko--Nikulin and the CHL denominator constructions give the sourced reflective presentations used for the corresponding rows \cite{GN97,GN98a,GN98b,GovindarajanKrishna2009,Govindarajan2011,GovindarajanSamanta2019}.  Every row has a minimal virtual product superspace obtained from its Fourier exponents, together with a Fock determinant and a diagonal Fredholm operator.  A sourced denominator presentation refines this product model by its Borcherds--Cartan bracket, Weyl group, and quantum supergroup.  The half-integral rows are formed on their multiplier covers.
\[
\begin{array}{c|c|c|c}
\text{form}&\text{group}&\text{weight}&\text{character order}\\ \hline
\Delta_5&\operatorname{Sp}_4(\Z)&5&2\\
\Delta_2&\Gamma_2&2&4\\
\Delta_1&\Gamma_3&1&6\\
\Delta_{1/2}&\Gamma_4&1/2&8\\
\nabla_3&\Gamma^{(2)}_0(2)&3&2\\
\nabla_2&\Gamma^{(2)}_0(3)&2&2\\
\nabla_{3/2}&\Gamma^{(2)}_0(4)&3/2&4\\
Q_1&\Gamma_2(2)&1&4
\end{array}
\]

\begin{theorem}[Eight-row algebraic realization]
Each diagonal-divisor row has a proper-height virtual root superspace, a super-Fock determinant, a colored abelian current theory, and a trace-class Fredholm denominator operator whose character is the corresponding automorphic product.  For every row equipped with its reflective denominator presentation, the finite-window Borcherds quiver, ordinary Hall category, three-Calabi--Yau completion, and parity-refined quantum supergroup refine this product model and recover the same denominator character.
\end{theorem}
\begin{proof}
Local finiteness of the product exponents gives the graded superspace and its proper-height completion.  The Fock and Fredholm constructions depend only on these signed multiplicities and reproduce the product coefficientwise.  A reflective denominator presentation supplies a symmetrizable Borcherds superdatum; the quiver, Hall, Ginzburg, current, and Hong constructions then apply to that datum and preserve its PBW denominator.
\end{proof}

\chapter{Algebraic and physical closure}
\begin{theorem}[Universal Igusa realization theorem]
The constructions have the following inputs and outputs. An arrow
labelled by a construction records its dependence on that input:
\[
\begin{gathered}
 \phiweak\xmapsto{\mathrm{Borch}}\Dmon,
 \qquad
 \phiweak\xmapsto{\mathrm{root\ datum}}\gIg,
 \qquad
 A_\Delta\xmapsto{\mathrm{quiver}}Q_{\Delta,H}
 \xmapsto{\mathrm{Hall}}\mathcal B_\Delta,\\
 kQ_{\Delta,H}\xmapsto{\Pi_3}\Gamma_{\Delta,H},
 \qquad
 \gIg\xmapsto{\Cur_E}\mathcal F^{\mathrm{hyb}}_{\Delta,E},
 \qquad
 \gIg\xmapsto{\mathrm{Hong}}\widehat U_q(\gIg),\\
 \gIg\xmapsto{\mathrm{root\ spectrum}}\mathbb A_Z,
 \qquad
 e^{-\rho}\operatorname{sPf}_F(\mathbb A_Z)=\DenIg(Z)=\Dmon(2Z),\\
 \mathcal Z^{\red}_{K3\times E}=-(\jmath^*\DenIg)^{-2}=-\ChiTen^{-1}.
\end{gathered}
\]
The system contains parity-refined one-particle spaces, an ordinary Hall extension product, the Borcherds root bracket, a local--wrapped collision law, a quantum supergroup, an automorphic orientation line, a first-order Fredholm determinant, and the reduced scalar square.
\end{theorem}
\begin{proof}
The Borcherds lift, Gritsenko--Nikulin presentation, and primitive
reduced curve-counting theorem give the scalar identities. The
enveloping and Chevalley maps are induced by the presentation.
The quiver Hall algebra and finite Calabi--Yau completion have their
respective quiver and path-algebra inputs. Their common grading
specifies the degrees a comparison must preserve. A map between
those constructions requires the additional functors, coefficient
maps, and pairings described above.
\end{proof}


\part{Technical appendices to the preceding book}
\chapter{Exact theta-series calculation}\label{app:theta-computation}
Set $t=q^{1/8}$ and $x=r^{1/2}$. Then
\[
\begin{aligned}
\vartheta_2(\tau,z)&=\sum_{n\in\Z}t^{4n^2+4n+1}x^{2n+1},\\
\vartheta_3(\tau,z)&=\sum_{n\in\Z}t^{4n^2}x^{2n},\\
\vartheta_4(\tau,z)&=\sum_{n\in\Z}(-1)^nt^{4n^2}x^{2n}.
\end{aligned}
\]
Squaring numerator and denominator, dividing in $\Q[x^{\pm1}][[t]]$, and summing the three quotients gives the coefficient table in Chapter~1. The calculation uses integer convolution and formal-series inversion; fractional powers cancel before conversion from $(t,x)$ to $(q,r)$.

\chapter{The root reflection matrices}
In the basis $(f_2,f_3,f_{-2})$, the form is
\[
 G=\begin{pmatrix}0&0&-1\\0&2&0\\-1&0&0\end{pmatrix}.
\]
The three reflection matrices are
\[
 S_1=\begin{pmatrix}1&4&4\\0&-1&-2\\0&0&1\end{pmatrix},\quad
 S_2=\begin{pmatrix}1&0&0\\-2&-1&0\\4&4&1\end{pmatrix},\quad
 S_3=\begin{pmatrix}1&0&0\\0&-1&0\\0&0&1\end{pmatrix}.
\]
Direct multiplication gives
\[
 S_i^2=I,
 \qquad S_i^{\mathsf T}GS_i=G.
\]

\chapter{Automorphic, homological, and Hall comparisons}
The half-index determinant, theta normalization, Borcherds denominator, reduced scalar theorem, charge Gram map, and geometric Hall comparison datum belong to the following constructions:
\begin{enumerate}
\item the automorphic square root is completed entirely within the Borcherds product theorem;
\item the BKM bracket comes from the Gritsenko--Nikulin algebra;
\item the bar--Chevalley comparison gives the denominator a homological source;
\item the current and BF constructions give an exact first-order algebraic model;
\item orientation data refine determinant lines, while an operator Pfaffian requires an actual Fredholm family;
\item Weyl reflections are attached to real roots, and imaginary roots enter through multiplicities;
\item the compact Hall comparison uses additive Mukai charge, the polarized central extension, and hybrid local--wrapped sectors.
\end{enumerate}
